\section*{第 \ref{ch:g9:climate-past-and-present} 章——地球的气候，过去与现在}

\begin{solution}{exo:g9:climate-past-and-present:1}
(a) 约 \qty{-19}{\degreeCelsius}。(b) 约 \qty{15}{\degreeCelsius}。
(c) $15 - (-19) = 34$：约 \qty{34}{\degreeCelsius}。(d) 任选两种：
二氧化碳、水蒸气、甲烷。
\end{solution}

\begin{solution}{exo:g9:climate-past-and-present:2}
(a) \omterm{def:g1:watching-the-weather:weather}{天气}。(b) \omterm{def:g5:weather-and-climate:climate}{气候}。(c) \omterm{def:g5:weather-and-climate:climate}{气候}（以及\omterm{def:g9:climate-past-and-present:climate-change}{气候变化}：
几十年间的变化）。(d) \omterm{def:g1:watching-the-weather:weather}{天气}。
\end{solution}

\begin{solution}{exo:g9:climate-past-and-present:3}
(a) \qty{19}{cm}，十一年。(b) $19 \div 11 \approx \qty{1.7}{cm}$。
(c) 每年的积雪都落在之前各年积雪之上，因此
层位越深，形成得越早。
\end{solution}

\begin{solution}{exo:g9:climate-past-and-present:4}
(a) $22.4 - 21.0 = +\qty{1.4}{\degreeCelsius}$。(b) $20.3 - 21.0 =
\qty{-0.7}{\degreeCelsius}$。(c) \qty{0}{\degreeCelsius}。(d) 该月
比 1991--2020 年同月的平均值更冷。
\end{solution}

\begin{solution}{exo:g9:climate-past-and-present:5}
(a) 从 1880 年代到 1930 年代、1950 年代和 1960 年代。(b) 约
\qty{0.25}{\degreeCelsius} 和 \qty{0.80}{\degreeCelsius}。(c) $0.80 -
0.25 = \qty{0.55}{\degreeCelsius}$。(d) 不对：从 1880 年代到
1910 年代各十年略有降温，1950 年代又比
1940 年代稍冷；但自 1970 年代起，每个十年
都比前一个十年更暖。
\end{solution}

\begin{solution}{exo:g9:climate-past-and-present:6}
(a) $(324.62 - 315.98) \div 10 \approx \qty{0.86}{ppm}$ 每年。(b)
$(427.35 - 401.01) \div 10 \approx \qty{2.63}{ppm}$ 每年。(c) $2.63
\div 0.86 \approx 3.1$：约快三倍。(d) $(427.35 - 315.98)
\div 315.98 \times 100 \approx 35~\%$。
\end{solution}

\begin{solution}{exo:g9:climate-past-and-present:7}
(a) $8 \times 10^5$ 年。(b) 约 \num{120000} 至 \num{130000} 年
前：此时温度曲线升至 0 以上。(c) 约 \qty{275}{ppm}。
(d) 是的：两条曲线的峰与谷出现在相同的时刻。
\end{solution}

\begin{solution}{exo:g9:climate-past-and-present:8}
(a) $3.2 \times 10 = \qty{32}{mm} = \qty{3.2}{cm}$。(b) $2100 - 2025 =
75$ 年；$75 \times 3.2 = \qty{240}{mm}$，约 \qty{24}{cm}（若上升加快，
还会更多）。(c) 海水受热膨胀；陆地上的冰
融化，其水流入大海。(d) 不会：\omterm{def:g8:ice-on-earth:sea-ice}{海冰}本就漂浮在海面上，
其融化不会使\omterm{def:g4:world-physical-map:altitude}{海平面}上升。
\end{solution}

\begin{solution}{exo:g9:climate-past-and-present:9}
(a) 90。(b) 66。(c) 99。(d) \emph{很可能}指成功的机会至少为 100 中的 66，
并非确定无疑：报纸使这一表述听起来比
科学家所说的更为肯定。
\end{solution}

\begin{solution}{exo:g9:climate-past-and-present:10}
某一地点的一个寒冷冬季是\emph{\omterm{def:g1:watching-the-weather:weather}{天气}}：它说的是
某处几个月内发生的情况。\omterm{def:g9:climate-past-and-present:global-warming}{全球变暖}则是
\emph{\omterm{def:g5:weather-and-climate:climate}{气候}}的变化，以几十年为尺度测量，并且是\emph{全球}的：它是
整个地球的平均值。某地的寒冬可以出现，
而整个地球总体上却比以前更暖，十年平均值
仍在持续上升。
\end{solution}

\begin{solution}{exo:g9:climate-past-and-present:11}
(a) $280 - 190 = \qty{90}{ppm}$，历时约 $2 \times 10^4$ 年。(b) $427
- 280 = \qty{147}{ppm}$，历时 $2025 - 1750 = 275$ 年。(c) 第二次
上升幅度更大（147 对 90~ppm）。速率：$90 \div (2 \times 10^4) =
4.5 \times 10^{-3}$~ppm 每年，而 $147 \div 275 \approx 5.3 \times
10^{-1}$~ppm 每年：约快 120 倍。(d) 2019 年的
浓度高于至少两百万年以来的任何时期（可信度高）。
\end{solution}

\begin{solution}{exo:g9:climate-past-and-present:12}
(a) 两者几乎相等：1.07 对 \qty{1.06}{\degreeCelsius}。最佳
估计认为，人类影响几乎解释了所观测到的全部
变暖。(b) 不是：自然因素合计使气温变化至多
\qty{0.1}{\degreeCelsius}，比所观测到的变暖小十倍。
(c) 人为导致的变暖至少有 100 中 66 的机会
介于 0.8 与 \qty{1.3}{\degreeCelsius} 之间。
\end{solution}

\begin{solution}{pb:g9:climate-past-and-present:1}
\textbf{1.} 因为它是参考时段：距平是
相对于该时段自身的平均值测量的。

\textbf{2.} 更冷，低 \qty{0.22}{\degreeCelsius}。

\textbf{3.} 更暖，高 \qty{0.98}{\degreeCelsius}。

\textbf{4.} $0.98 - (-0.22) = \qty{1.20}{\degreeCelsius}$。

\textbf{5.} $427 - 316 = \qty{111}{ppm}$。

\textbf{6.} $111 \div 316 \times 100 \approx 35~\%$。

\textbf{7.} $427 - 300 = \qty{127}{ppm}$。

\textbf{8.} 二氧化碳吸收暖的地表所放出的部分红外线，
并将其中一部分送回地面。空气中二氧化碳越多，
回到地表的能量就越多，地表随之变暖。

\textbf{9.} $2025 - 1993 = 32$ 年；$32 \times 3.2 \approx
\qty{102}{mm}$，约 \qty{10}{cm}。

\textbf{10.} 变暖的海水膨胀；陆地上的冰（\omterm{def:g8:ice-on-earth:glacier}{冰川}、\omterm{def:g8:ice-on-earth:ice-sheet}{冰盖}）
融化并流入大海。

\textbf{11.} 是的，大小相同：约 \qty{1.2}{\degreeCelsius} 对
\qty{1.09}{\degreeCelsius}。两者略有差异，是因为所取时段不同
（1880--1899 年而非 1850--1900 年；2015--2024 年而非
2011--2020 年：较近的年份更暖），并且因为 IPCC
综合了多项记录，而不只是 NASA 的记录。

\textbf{12.} “自 19 世纪末以来，地球表面已
变暖约 \textbf{\qty{1.2}{\degreeCelsius}}。”
\end{solution}
